\chapter{Getting started}
\label{ch:start}

\section{Installation}

The build is that of GROMACS~4.5.4 with CMake.
FFTW~3 (single precision) is required; it is also what the electron model uses for its grid convolutions.
\begin{lstlisting}
cd gromacs-4.5.4-MCMD
mkdir build && cd build
cmake .. -DCMAKE_INSTALL_PREFIX=/path/to/install
make -j8 install
\end{lstlisting}
Use the \code{grompp} of this build to write run inputs: the \file{.tpr} format carries the \opt{mcmd-*} options
(format 80), which a stock GROMACS cannot read (and an older stock GROMACS tool would silently misread).
Set \env{GMXLIB} to the build's \file{share/top} if the force fields are not found.

\section{The example}

\file{example/} holds a complete run: hen egg-white lysozyme (PDB 1AKI, 1960 atoms) in vacuum, exploded by a
Gaussian pulse of $10^{11}$ photons at 1000\,eV, 2.94\,fs FWHM, focused to 100\,nm.
\begin{lstlisting}
cd example
pdb2gmx -f 1aki.pdb -ff charmm27 -water none
grompp  -f exp.mdp -c conf.gro -p topol.top -o explode.tpr
mdrun   -deffnm explode -v -nt 4 -pd -ionize
\end{lstlisting}
(\file{run\_example.sh} does the same.)
The run is 15\,000 steps of 2\,as, i.e.\ 30\,fs with the pulse peaking at $t=0$, and takes about 10 minutes on 4
threads.
No atomic data files are needed: the module computes them for 1000\,eV at start-up (Section~\ref{sec:builtin}).
The model is the default one; the ``Model'' block of \file{exp.mdp} lists its main options.

\paragraph{What to expect.}
\opt{mcmd-seed = -1} draws a new random seed every run, so the numbers differ by about a percent between runs.
Roughly:
\begin{itemize}
  \item every atom ends up ionized; the mean charge is about 3.5 per atom (hydrogen fully stripped, carbon $\sim$5.4,
    nitrogen $\sim$6.2, oxygen $\sim$7.1, sulfur $\sim$13.6);
  \item about 4000 photoionizations, 800 Auger decays, 2100 electron-impact ionizations and 820 charge transfers;
  \item of the $\sim$6900 electrons released, $\sim$1350 have escaped and $\sim$5600 are still trapped at 30\,fs, at a
    temperature of $\sim$18\,eV; the radius of gyration has grown 3.6-fold, to $\sim$5.7\,nm.
\end{itemize}
\code{tools/mcmd\_output.py plot} draws an overview of the whole run (Section~\ref{sec:viewer}).

\section{The run summary}
\label{sec:summary}

At start-up \code{mdrun} prints a summary of what the module will do, and writes the same block to \file{md.log}
(after GROMACS's own parameter listing): the system, the pulse, the run length and seed, the atomic data found, every
process and whether it is on, and the output.
For the example:
\begin{lstlisting}
==================================================================================
 MolDStruct MC/MD
 System        1960 atoms (H 959, C 613, N 193, O 185, S 10), all neutral at the start
 Pulse         1000 eV, Gaussian, FWHM 2.9435 fs, peak at 0 fs, 1e+11 photons into 100 nm
               (1.27e+13 photons/um^2); bandwidth 0.003
 Run           15000 steps x 2 as = 30 fs; seed 3104721958 (chosen at start, mcmd-seed = -1); 4 ranks
 Atomic data   built in: screened-hydrogenic shells (K, L, M) at 1000 eV, mcmd-atomic-data = hydrogenic
               (written to ./simulation_output/Atomic_data)

 Processes
   photoionization, Auger, fluorescence     on (always)
   charge transfer                          on   hydrogen donors only, H-H off, downhill off
   trapped electrons (Poisson-Boltzmann)    on   mcmd-electrons = pbgrid
       grid                                 s_e 0.25 nm, 2.5 cells/s_e, box <= 3 R, <= 384 cells,
                                            solve every 0.2 fs, Newton tol 0.01
       evaporation                          on
   ion-ion forces                           all pairs
   electron-impact ionization               on   Lotz, built in (mcmd-coll-data = lotz)
   three-body recombination                 off
   collisional excitation                   off
   field ionization                         off

 Output        simulation_output/: transitions.bin (mcmd-detailed-output = 1), charges_over_time.bin every 500 steps;
               autostop off
 Settings      all at the validated production values
==================================================================================
\end{lstlisting}
Settings outside the validated model are listed as warnings instead of the last line --- an MD step above 2\,as, a
PB solver coarser than the converged settings, a non-default smoothing, three-body recombination, collisional
excitation, H--H charge transfer, all pairs above 10\,000 atoms, harmonic bonds that cannot break, a pulse without
photons.
The run goes on; the warnings say that the result is not covered by the validation of
Appendix~\ref{app:validation}.

\section{Running your own system}

\begin{enumerate}
  \item Prepare the system as for any vacuum GROMACS run: \code{pbc = no}, no pressure or temperature coupling during
    the explosion.
    Equilibrate it beforehand without \code{-ionize}.
  \item Set \code{morse = yes} so that \code{grompp} turns the bonds into Morse potentials, which can break
    (Chapter~\ref{ch:ff}).
  \item Use \code{dt = 0.000002} (2\,as) and the pulse options \opt{mcmd-pulse-*} (Chapter~\ref{ch:mdp}); note that
    the pulse times are in fs.
  \item The atomic data are computed for the photon energy of the pulse (H to Ar except He); for other data supply
    \file{Atomic\_data/} and set \opt{mcmd-atomic-data = files} (Chapter~\ref{ch:data}).
  \item Run \code{mdrun -ionize}, with \code{-nt N -pd} for more than one thread (Appendix~\ref{app:parallel}).
\end{enumerate}
Ion-ion Coulomb interactions are taken care of by \opt{mcmd-pb-long-range = auto}: all pairs up to 10\,000 atoms,
a 2\,nm cutoff plus the long-range field from the grid above (\code{grompp} sets the cut-offs and says so).

\section{Reproducible runs and continuation}
\label{sec:seed}

The Monte Carlo draws come from a random seed.
With \opt{mcmd-seed = -1} (the default) it is drawn from the clock and the process id at start-up, so every run is
a different realization --- which is what an ensemble needs.
The seed is printed in the summary; set \opt{mcmd-seed} to it to repeat a run exactly (on the same number of
ranks, Appendix~\ref{app:parallel}).
Give every member of an ensemble its own seed, or leave it at $-1$.

Runs can be continued from a GROMACS checkpoint with \code{mdrun -cpi}.
At every checkpoint the module writes \file{simulation\_output/mcionize\_state.bin} (the previous one kept as
\file{mcionize\_state\_prev.bin}): the electronic configurations, the electron cloud's state, the random streams, the
reference radius the expansion is measured against, and the length of every output file.
On restart the state of the checkpoint's step is read back, the output files are cut back to that step and appended
to.
The state file is required: the reference radius in particular cannot be recovered from an expanded sample.
To extend a finished run, raise \code{nsteps} with \code{tpbconv -extend} and continue with \code{-cpi}.

\section{The regression test}

\file{test/regression.sh} runs the example for 300 steps with a pinned seed in two cases --- \code{none} (no
electron cloud) and \code{full} (the defaults, on a coarse grid) --- and checks that the event counts match the
stored references exactly and that the charge is conserved.
Then, with \opt{mcmd-atomic-data = files} on the tables the full case wrote, it checks that they reproduce the run
bit for bit, that a run without \code{-ionize} writes no ionization output, that the total energies are required
exactly when charge transfer or collisions need them, and that the run is unchanged without the statistical weights.
It takes under a minute; \code{--update} regenerates the references after a deliberate change of the physics.
The references are only meaningful for the build and machine that wrote them.
